\section{Telemetry 与 Spend Control：可见不等于可控}

上一章说明 consumption economics 必须落实到 operation-level attribution；但“能计量”还没有回答异常发生时谁来行动。本章把近实时 telemetry 与 spend controls 连成闭环：前者把部署变化映射到 latency、error、resource 和费用，后者决定越过阈值后通知、触发 webhook、降级还是暂停。核心不是拥有更多 dashboard，而是用可解释信号做出与业务价值相称的控制决策。

\subsection{把 deploy、latency、error 和 spend 放在同一时间线}

本节先建立可观测性的因果顺序：一次 deployment 改变了代码与配置，随后才可能改变性能、错误率、资源消耗和账单。读图时不要把四条独立曲线并排观看，而要沿同一版本时间线追问“哪次变化造成了什么后果”，并区分业务流量增长、应用回归与平台计量变化。

\lecturefigure{11-telemetry-loop.png}{Telemetry 应把版本变化与资源、性能和费用后果关联起来。}{本地字幕 22:50--23:45；概念重绘。}

\begin{knowledgebox}{读图：为什么 telemetry 自身也是昂贵基础设施}
高基数 labels、分钟级账单、跨区域日志和长保留期都会增加存储与查询成本。平台收费的一部分确实是在维护这套反馈基础设施。用户应明确哪些信号需要实时、哪些可采样、哪些必须长期保留用于审计。
\end{knowledgebox}

\subsection{Soft cap 与 hard cap 的目标不同}

Soft cap 通过通知或 webhook 促使人和自动化采取行动，但流量继续服务；hard cap 在阈值后暂停，限制最坏支出，却可能切断正在产生收入的业务。当前 Vercel Spend Management 文档还明确指出检查并非连续，暂停可能在越过阈值数分钟后发生，因此 hard cap 也不是精确断路器。

\lecturefigure{12-spend-controls.png}{Soft control 保护学习与连续性，hard control 保护最坏支出。}{本地字幕 23:40--24:50；Vercel Spend Management 官方文档。}

\begin{importantbox}{预算阈值应来自业务函数}
可把 hard-stop 决策粗略写为：当继续服务的预期损失
\[
L_{continue}=P(\text{abuse})\cdot C_{abuse}+C_{runaway}
\]
超过暂停造成的
\[
L_{pause}=R_{lost}+C_{recovery}+C_{trust}
\]
时才暂停。这里 $R_{lost}$ 是丢失收入，$C_{trust}$ 是用户信任损失。实际系统需要多级限流、attack detection 和人工升级，而不是只设一个金额。
\end{importantbox}

\begin{warningbox}{Hard cap 不是实时安全边界}
计量、汇总、阈值判断和暂停传播都存在延迟。对高风险 AI endpoint，还应使用 request rate limit、per-user quota、model budget、circuit breaker 和 abuse monitoring，不能只依赖月底预算。
\end{warningbox}

\teachervoice{课堂案例中，客户因 hard cap 在高流量时停止服务而损失收入。无论数字是否精确，结论很清楚：成本控制必须理解业务价值。Soft cap 用于学习，hard cap 用于不可接受风险；两者不能互换。}

\subsection{本章小结}

Telemetry 提供因果线索，Spend Control 提供行动。有效系统要把性能、成本、业务价值和执行延迟共同纳入阈值设计。

